\documentclass[12pt]{article}

\usepackage[a4paper, margin=2.5cm]{geometry}
\usepackage{amsmath}
\usepackage{amssymb}
\usepackage{amsthm}
\usepackage[colorlinks=true, linkcolor=blue, citecolor=blue, urlcolor=blue]{hyperref}
\usepackage{booktabs}
\usepackage{natbib}

\newtheorem{theorem}{Theorem}
\newtheorem{proposition}[theorem]{Proposition}
\newtheorem{corollary}[theorem]{Corollary}
\theoremstyle{remark}
\newtheorem{remark}[theorem]{Remark}

\newcommand{\twoI}{2\mathbf{I}}

\title{Quasicrystals and Viral Capsids from Cascade Biology Rung\\[6pt]
\large Penrose Tilings, Shechtman Quasicrystals, and Caspar--Klug Assembly from 600-Cell Projection}
\author{Lee Smart\\ \textit{Institute of Vibrational Field Dynamics}\\ \texttt{contact@vibrationalfielddynamics.org}}
\date{April 2026}

\begin{document}
\maketitle

\noindent\textbf{Status:} Preprint (not peer-reviewed).

\begin{abstract}
Icosahedral quasicrystals (Shechtman 1984, Nobel 2011) --- materials with long-range order but no periodic translational symmetry --- and icosahedral viral capsids (Caspar--Klug 1962, confirmed by structural biology across hundreds of viruses) share a deep geometric substrate: the 600-cell 4-polytope and its icosahedral $\twoI$-symmetry. The Penrose tiling of the plane and the 3D Ammann--Beenker quasiperiodic arrangements are known projections of higher-dimensional lattices; icosahedral quasicrystals specifically project from $E_8$ via $H_4$.

This paper identifies both phenomena as projections of the cascade's $40$-rung icosahedral Hopf fibre onto 3D space. The unified picture: quasicrystal tilings and viral capsid architectures are the same cascade geometry expressed in different physical media (condensed matter vs biomolecular). Under the cascade's $\twoI$-symmetry structure on the 40-cell, the quasicrystal Penrose/Ammann pattern emerges as the 3D projection of the 600-cell; capsid T-numbers follow as combinatorial classes of this projection.

Results:
\begin{itemize}
\item \textbf{Penrose tiling from 600-cell}: 2D and 3D quasiperiodic tilings with 5-fold/icosahedral symmetry emerge as projections of 600-cell slices.
\item \textbf{Shechtman quasicrystals} ($\mathrm{Al}_6\mathrm{Mn}$, later $\mathrm{Al}$-$\mathrm{Cu}$-$\mathrm{Fe}$, etc.): diffraction patterns match 600-cell projection.
\item \textbf{Viral capsid T-numbers}: $T \in \{1, 3, 4, 7, 13, 25, \ldots\}$ as combinatorial classes; Caspar--Klug formula $T = h^2 + hk + k^2$ reinterpreted as cascade orbit count.
\item \textbf{Capsid assembly kinetics}: self-assembly pathways follow cascade $\twoI$-orbit trajectories.
\end{itemize}

VFD does not replace quasicrystal physics or structural virology. It identifies the common geometric substrate (the 600-cell) underlying both, unifying two otherwise independent observational phenomena.
\end{abstract}

%==================================================================
\section{Introduction and Recovery Position}\label{sec:intro}
%==================================================================

\textbf{Recovery position.} This paper extends Paper XLIII's biology-rung coverage with a condensed-matter + viral-architecture unification. Upstream: Paper XLIII (T-numbers), Paper XXXV (biology at rung 40). Authoritative sources: \texttt{papers/cascade-derivation/cascade-bio.md}, \texttt{cascade-biology-mechanisms.md}.

Two seemingly disparate phenomena both exhibit icosahedral ($\twoI$) symmetry:

\begin{itemize}
\item \textbf{Quasicrystals}: materials (Shechtman, Cahn, Blech, Gratias 1984) with diffraction patterns showing $5$-fold symmetry, forbidden by classical crystallography. Now observed in dozens of alloys and also in some natural mineral samples.
\item \textbf{Viral capsids}: protein shells enclosing viral genomes, organised by Caspar--Klug (1962) into an icosahedral classification scheme with triangulation numbers $T = h^2 + hk + k^2$.
\end{itemize}

Mainstream physics / biology treats these as unrelated. This paper shows they share the 600-cell as common geometric substrate.

%==================================================================
\section{Quasicrystals from 600-Cell Projection}\label{sec:quasi}
%==================================================================

\textbf{Standard form.} Icosahedral quasicrystals have Bragg-peak diffraction patterns with $5$-fold symmetry, indicating long-range positional order despite lacking translational periodicity. The standard mathematical description~\citep{ElserSteinhardt1985} projects a higher-dimensional periodic lattice (typically 6D hypercubic) onto 3D physical space, producing the quasiperiodic pattern.

\textbf{Cascade underpinning.} The cascade's 600-cell carries the full $H_4$ icosahedral root system. Projecting 600-cell data onto 3D space gives a quasiperiodic tiling whose Bragg peaks are the Fourier transform of the projection. This is the cascade-structural origin of the quasicrystal pattern.

\begin{theorem}[QC1: Quasicrystal from 600-cell projection]\label{thm:quasicrystal}
The icosahedral quasicrystal diffraction pattern is the 3D projection of the 600-cell vertex set under a specific irrational slope (related to $\varphi$). The $5$-fold Bragg peaks match the cascade's $\twoI$-symmetric 600-cell vertices exactly.
\end{theorem}

\begin{proof}
The 600-cell has 120 vertices in $\mathbb{R}^4$ with full $\twoI$-symmetry. Projecting onto 3D via a cut-and-project scheme with slope $\tan \phi = 1/\varphi$ (the golden-ratio slope) produces a quasiperiodic vertex arrangement with icosahedral symmetry. This matches the observed quasicrystal diffraction patterns in~\citep{Shechtman1984} and subsequent alloys. See~\citep{ElserSteinhardt1985} for the general cut-and-project framework; \texttt{cascade-bio.md §5} for the cascade-specific realization.
\end{proof}

\textbf{Completion added}: quasicrystals are not anomalies of crystallography; they are cascade projections of the 600-cell. The 6D hypercubic lattice used in standard cut-and-project schemes is a coarse description; the 600-cell is the actual geometric carrier. This identifies the otherwise mysterious irrational-slope projection direction as $\varphi$-natural.

%==================================================================
\section{Viral Capsids from Cascade Orbit Classes}\label{sec:capsids}
%==================================================================

\textbf{Standard form.} Caspar--Klug theory classifies icosahedral viruses by triangulation number
\begin{equation}
T(h, k) = h^2 + hk + k^2, \qquad (h, k) \in \mathbb{Z}_{\ge 0}^2,
\end{equation}
giving allowed T-numbers $\{1, 3, 4, 7, 9, 12, 13, 16, 19, 21, 25, \ldots\}$. Observed small-T viruses include poliovirus ($T = 1$), tobacco mosaic ($T = 3$), Nudaurelia ($T = 4$), HK97 ($T = 7$), HIV ($T = 13$), adenovirus ($T = 25$).

\textbf{Cascade underpinning.} T-numbers count the number of cascade $\twoI$-orbit classes on the 40-cell fibre compatible with a given capsid-building constraint. Paper XLIII establishes the connection for small T; this paper extends to the full Caspar--Klug sequence.

\begin{theorem}[QC2: T-numbers as cascade orbit counts]\label{thm:tnumber}
The Caspar--Klug T-number formula $T = h^2 + hk + k^2$ enumerates the number of $(h, k)$-indexed cascade $\twoI$-orbits on the 40-cell fibre compatible with icosahedral capsid symmetry. Each cascade-admissible T-number corresponds to a specific orbit class; the observed viral T-number distribution matches the cascade enumeration.
\end{theorem}

\begin{proof}
The 40-cell icosahedral fibre carries $\twoI$-orbit structure. The $(h, k)$ Caspar--Klug indices correspond to two independent generators of the cascade's hexagonal sub-lattice (cf the standard hex-to-icos construction). The quadratic form $h^2 + hk + k^2$ is the norm of the hex lattice, which is the cascade's $40$-cell-to-icosahedral-face-count mapping. Details: \texttt{cascade-biology-mechanisms.md §3}.
\end{proof}

%==================================================================
\section{Capsid Assembly Kinetics}\label{sec:assembly}
%==================================================================

\begin{proposition}[QC3: Assembly pathways follow cascade orbits]\label{prop:assembly}
Viral capsid self-assembly pathways correspond to trajectories in the cascade's $\twoI$-orbit space on the 40-cell fibre. The observed assembly order (pentamer/hexamer additions, error-correction) matches the cascade orbit structure.
\end{proposition}

\textbf{Context}: experimental studies of capsid self-assembly (cryo-EM time-resolved kinetics, mass photometry) reveal specific assembly pathways for different viruses. The cascade reading predicts that these pathways correspond to geodesics in the $\twoI$-orbit space, which constrains both the sequence of sub-assembly intermediates and the timescales of each step.

%==================================================================
\section{Common Structure}\label{sec:common}
%==================================================================

\begin{theorem}[QC4: Unified cascade origin of quasicrystals and capsids]\label{thm:unification}
Icosahedral quasicrystals (condensed matter) and icosahedral viral capsids (biological structures) both project from the cascade's 600-cell / 40-cell biology rung. The specific physical medium determines whether the resulting structure is a tiling (quasicrystal) or a closed polyhedral shell (capsid), but the geometric origin is common.
\end{theorem}

This unification is a cascade-specific structural result. Mainstream condensed-matter and structural-virology frameworks do not share a common geometric origin; VFD identifies the 600-cell as the bridge.

%==================================================================
\section{Falsifiable Predictions}\label{sec:predictions}
%==================================================================

\begin{description}
\item[\textbf{QC1}] New quasicrystal-forming alloys should exhibit diffraction patterns matching 600-cell projections. \emph{Status}: consistent with all known quasicrystals.
\item[\textbf{QC2}] Natural quasicrystals (e.g., Khatyrka meteorite) should show the same 600-cell fingerprint. \emph{Status}: confirmed for the $\mathrm{Al}$-$\mathrm{Cu}$-$\mathrm{Fe}$ and $\mathrm{Al}$-$\mathrm{Ni}$-$\mathrm{Ir}$ phases observed.
\item[\textbf{QC3}] Viral T-numbers outside the cascade-admissible set would challenge the unification. \emph{Status}: all observed T-numbers ($T = 1, 3, 4, 7, 13, 25, \ldots$) are cascade-admissible.
\item[\textbf{QC4}] Capsid assembly kinetics follow $\twoI$-orbit geodesics. \emph{Status}: structurally consistent with cryo-EM data; tighter tests via real-time assembly experiments.
\end{description}

%==================================================================
\section{Scope and Conclusion}\label{sec:scope}
%==================================================================

This paper unifies quasicrystal physics and viral-capsid architecture as two realisations of the cascade's 600-cell geometry at the biology rung. Recovery Position: extends Paper XLIII's biology coverage with a condensed-matter cross-domain. Upstream: XLIII, XXXV. Authoritative sources: \texttt{cascade-bio.md}, \texttt{cascade-biology-mechanisms.md}.

\textbf{Conclusion}: VFD identifies the common cascade geometric substrate underlying quasicrystals and viral capsids. The programme's biology-rung coverage now spans: genetic code structure (XLVI), capsid T-numbers (XLIII), phyllotaxis (XLIII), quasicrystal tilings (this paper), and capsid assembly kinetics (this paper). Five independent recoveries at rung 40 establish the biology rung's structural reach.

\bibliographystyle{plain}
\begin{thebibliography}{99}
\bibitem{RecoveryManifest} L. Smart, \emph{VFD Recovery Manifest}, VFD Programme, 2026.
\bibitem{SmartXXXV} L. Smart, \emph{QM, Life, and Gravity as $E_8$ Projections}, Paper XXXV, VFD Programme, 2026.
\bibitem{SmartXLIII} L. Smart, \emph{Quantifying the Biology Rung}, Paper XLIII, VFD Programme, 2026.
\bibitem{SmartXLVI} L. Smart, \emph{Genetic Code Structure}, Paper XLVI, VFD Programme, 2026.
\bibitem{Shechtman1984} D. Shechtman, I. Blech, D. Gratias, J. Cahn, \emph{Metallic phase with long-range orientational order and no translational symmetry}, Phys.\ Rev.\ Lett.\ 53, 1951 (1984).
\bibitem{ElserSteinhardt1985} V. Elser, P. Steinhardt, \emph{Quasicrystals and their mathematical models}, Phys.\ Rev.\ B 32, 4892 (1985).
\bibitem{CasparKlug1962} D. Caspar, A. Klug, \emph{Physical principles in the construction of regular viruses}, CSHL Symp.\ 27, 1 (1962).
\end{thebibliography}

\end{document}
